\section{Theoretical Background}

\subsection{Information Security Risk Management}
Information security risk management is a structured process used to identify, assess, and treat risks that may affect the confidentiality, integrity, and availability of information assets. It provides organizations with a systematic approach for understanding potential threats, evaluating vulnerabilities, and prioritizing mitigation efforts based on the likelihood and potential impact of adverse events.

In contrast to purely technical security measures, risk management emphasizes decision-making under uncertainty. Rather than attempting to eliminate all risks, organizations seek to reduce risk to an acceptable level through a combination of preventive, detective, and corrective controls. This process typically involves identifying assets, analyzing threats and vulnerabilities, estimating risk levels, and selecting appropriate treatment strategies such as risk mitigation, transfer, avoidance, or acceptance.

\subsection{Risk Management Frameworks and Standards}
Several internationally recognized frameworks and standards provide structured guidance for conducting information security risk management. While they may differ in terminology and emphasis, they share a common objective; to provide systematic methods for identifying, analysing, and treating risk.

ISO27005 offers guidance specifically focused on information security risk management within the context of an Information Security Management System (ISMS). It emphasizes a cyclical process involving context establishemnt, risk assessment, risk treatment, risk acceptance, and communication.

The NIST Risk Management Framework provides a structured, lifecycle-oriented approach integrating security and risk management activities into system development and operation. It focuses on categorization, control selection, implementation, assessment, authorization, and continuous monitoring.

OCTAVE (Operationally Critical Threat, Asset, and Vulnerability Evaluation) focuses on organizational risk evaluation driven by internal knowledge and operational context. It places strong focus on identifying critical assets and assessing risks from a business perspective.

FAIR (Factor Analysis of Information Risk) differs from qualitative approaches by providing a model for quantitative risk analysis. It aims to estimate risk in financial terms by decomposing risk into measurable factors such as threat event frequency and loss magnitude.

In addition, guidance from ENISA and national authorities such as NSM provides recommendations for governance, maturity, and practical implementation within European and Norwegian contexts.

Although these frameworks vary in structure and depth, they collectively define expectations regarding documentation, traceability, and repeatability in risk management processes. These expectations form an important foundation for evaluating tool support in later chapters.

\subsection{Role of Tool Support in Risk Management}
While risk management frameworks define processes and principles, practical implementation often relies on tool support. Tools can assist in structuring workflows, maintaining documentation, tracking risk treatment measures, and generating reports for internal governance or regulatory compliance purposes.

At a basic level, organizations may use spreadsheets or manual documentation systems to manage risk registers. More advanced solutions provide structured databases, predefined threat libraries, automated risk calculations, integration with asset inventories, and workflow management capabilities. Enterprise-level platforms may additionally integrate with governance, risk, and compliance (GRC) systems, security monitoring tools, or ticketing systems.

Tool support enhance consistency, scalability, and auditability, particularly in larger organizations where risk management processes involve multiple stakeholders and recurring assessments. Tools do however not replace methodological understanding. Their effectiveness depends on how well they align with established frameworks and organizational needs.

For this reason, evaluation risk management tools requires examining not only functional capabilities but also their alignment with recognized standards, support for traceability and documentation, and suitability for different operational contexts. These considerations provide the foundation for the comparative evaluation conducted in subsequent chapters.